\section{Provider Configuration and Execution Accounting}
\label{supp:provider}

\paragraph{Provider policy.}
Phase~4J used Alibaba \texttt{qwen3.7-max} in the pinned Singapore workspace,
with thinking disabled, temperature unset, a 90-second timeout, 4,096 maximum
output tokens, zero retries, and no schema repair.  The provider required
usage metadata, allowed no mock fallback, and executed at concurrency one.
Phase~5B used the same provider/model and region; its frozen record specifies
temperature 0.0.  Hosted services can remain nondeterministic under settings
intended to be deterministic \citep{AtilEtAl2025Nondeterminism}.

\paragraph{Ablation budgets.}
A and B allow 4,096 output tokens for each of two calls; C and D allow 2,048
for each of four calls.  Nominal condition-level capacity is therefore 8,192
tokens throughout.  The frozen scheduler rotates condition order by cell and
records 216 planned calls.  Retries and repairs are zero.

\paragraph{Original-artifact-backed accounting.}
The recovered aggregate and sentinel support 72 intended units, 70 completed
units, two failed units, elapsed time 2,396.74 seconds, and 216 planned calls.
The historical field named total physical dispatches records the plan and is
not evidence that all 216 requests were sent.

\paragraph{Transcript-backed accounting.}
The execution transcript reports 214 attempted calls, 212 responses, 488,640
input tokens, 71,820 output tokens, and estimated cost \$0.2816.  Because the
raw request/response ledger did not survive in the recovered run directory,
these quantities remain transcript-backed and are not canonical scientific
evidence.

\paragraph{Authorization lifecycle.}
Both live experiments were bound to a freeze, provider/model policy,
confirmation phrase, sentinel, and one-execution maximum before execution.
The retained Phase~5B authorization has status
\texttt{consumed-live-execution-completed}.  Authorization records carry no
API key.
